%----------------------------------------------------------------------------------------
% TRANSFORMATIES VAN GRAFIEKEN. Alle opgaven zijn eigen werk.
%
% Opbouw: Probleem, Start, leerstof (theorie, werkwijze of uitleg, voorbeeld, korte oefening), en op een
% nieuwe pagina de Opdrachten (door elkaar, stijgende moeilijkheid), Koppel en Test jezelf.
% Labels van dit hoofdstuk beginnen met h2-.
%----------------------------------------------------------------------------------------

\mychapter{2}{Transformaties van grafieken}

\begin{probleem}
\vspace{-1mm}\naast{De grafiek toont de temperatuur in Brussel op een winterdag, in graden Celsius, $t$ uur na middernacht. Je vriendin in New York wil dezelfde metingen zien zoals ze dat gewend is: met de temperatuur in graden Fahrenheit, en met de klok van New York, waar het 6 uur vroeger is dan in Brussel. Een temperatuur van $T$\,$^\circ$C is $1{,}8\,T + 32$\,$^\circ$F. Hoe ziet haar grafiek eruit? Om hoe laat was het volgens haar klok het warmst, en hoeveel graden Fahrenheit was het toen? Waar zie je op haar grafiek dat het vroor?}{\begin{tikzpicture}[x=.24cm,y=.36cm,font=\scriptsize]
  \draw[black!14,very thin,xstep=.72cm,ystep=.36cm] (0,-3) grid (24,7);
  \draw[->,grijs] (0,0) -- (25.5,0) node[right,text=black] {$t$};
  \draw[->,grijs] (0,-3) -- (0,7.6) node[above,text=black] {$^\circ$C};
  \foreach \t in {6,12,18,24}{\node[below,inner sep=1.5pt] at (\t,0) {$\t$};}
  \foreach \t in {-2,2,4,6}{\node[left,inner sep=1.5pt] at (0,\t) {$\t$};}
  \kromme{0}{24}{-4*cos(15*(\x-3))+2}
\end{tikzpicture}}
\end{probleem}

\lkr{Leerplan \pakketnaam: LPD \alleenvoor{4}{9-11}\alleenvoor{1}{5-6}\alleenvoor{2,3}{9 en 14 (de transformaties bij de exponentiële functies en bij de algemene sinusfunctie)}\alleenvoor{4}{ (algemene functies)}\alleenvoor{1}{, met de transformaties die LPD 9 (exponentiële functies) en LPD 12 (sinusfuncties) gebruiken}, het deel over transformaties van grafieken. In de cursussen over exponentiële, goniometrische en rationale functies komen dezelfde transformaties terug: daar volstaat dan een verwijzing naar dit hoofdstuk. Laat leerlingen telkens één bijzonder punt volgen (een top, een beginpunt): dat maakt de transformatie zichtbaar en controleerbaar.\alleenvoor{1}{ In Basis zonder bewijzen: de uitleg bij elke regel volstaat. Spiegelen om de $y$-as blijft, omdat Exponentiële functies er de grafiek van $\bigl(\frac12\bigr)^x$ mee uitlegt. Met schuifbalken in GeoGebra zien leerlingen wat elke parameter doet (wenk bij LPD 9 en 12).}\alleenvoor{2}{ In B+S zonder bewijzen: de uitleg bij elke regel volstaat (bewijsvoering is bij LPD 4 alleen een extra). Met schuifbalken in GeoGebra zien leerlingen wat elke parameter doet (wenk bij LPD 9 en 14). De verschoven $\sqrt{x}$ en $\frac1x$ bereiden de elementaire irrationale en homografische functies van LPD 17 voor.}\alleenvoor{3}{ In B+S' met bewijzen (bewijsvoering hoort bij LPD 4). Met schuifbalken in GeoGebra zien leerlingen wat elke parameter doet (wenk bij LPD 9 en 14). De verschoven $\sqrt{x}$ en $\frac1x$ bereiden de elementaire irrationale en homografische functies van LPD 18 voor.}}

Je vriendin kan elk punt van de grafiek opnieuw uitrekenen, maar dat hoeft niet. Elke temperatuur verandert op dezelfde manier (maal $1{,}8$, plus $32$), en elk tijdstip ook (6 uur vroeger). De grafiek van New York is dus de grafiek van Brussel in een andere vorm: verschoven en uitgerekt. Zo zijn veel functies een bekende functie in een andere vorm: verschoven, uitgerekt of gespiegeld. Wie de grafieken van een handvol basisfuncties kent, kan zo heel wat grafieken schetsen zonder tabel. Het idee is telkens hetzelfde: een grafiek is een verzameling punten (hoofdstuk 1), en een kleine verandering in het voorschrift verplaatst elk punt op dezelfde manier. We zoeken dus uit waar elk punt $(x, y)$ terechtkomt. Eén ding verrast daarbij: wat je \textit{binnen} de functie verandert, werkt omgekeerd.

\begin{oefeningen}[Start]
\oef{h2-start-1}

\oef{h2-verken}
\end{oefeningen}

%========================================================================================
\section{Basisfuncties}

Deze grafieken ken je best uit het hoofd, met hun domein en bereik.

\figuren{%
\begin{tikzpicture}[x=.36cm,y=.36cm,font=\scriptsize]
  \roosterleeg{-3}{1}{3}{-3}{1}{3}
  \kromme{-1.73}{1.73}{(\x)^2}
  \node[below,text=black] at (0,-3.2) {$x^2$};
\end{tikzpicture}\hspace{3mm}
\begin{tikzpicture}[x=.36cm,y=.36cm,font=\scriptsize]
  \roosterleeg{-3}{1}{3}{-3}{1}{3}
  \kromme{-1.44}{1.44}{(\x)^3}
  \node[below,text=black] at (0,-3.2) {$x^3$};
\end{tikzpicture}\hspace{3mm}
\begin{tikzpicture}[x=.36cm,y=.36cm,font=\scriptsize]
  \roosterleeg{-3}{1}{3}{-3}{1}{3}
  \kromme{0}{3}{sqrt(\x)}
  \vol{(0,0)}
  \node[below,text=black] at (0,-3.2) {$\sqrt{x}$};
\end{tikzpicture}\hspace{3mm}
\begin{tikzpicture}[x=.36cm,y=.36cm,font=\scriptsize]
  \roosterleeg{-3}{1}{3}{-3}{1}{3}
  \draw[accent,thick] (-3,3) -- (0,0) -- (3,3);
  \node[below,text=black] at (0,-3.2) {$|x|$};
\end{tikzpicture}\hspace{3mm}
\begin{tikzpicture}[x=.36cm,y=.36cm,font=\scriptsize]
  \roosterleeg{-3}{1}{3}{-3}{1}{3}
  \kromme{.34}{3}{1/\x} \kromme{-3}{-.34}{1/\x}
  \node[below,text=black] at (0,-3.2) {$\frac1x$};
\end{tikzpicture}}

\begin{theorie}
\renewcommand{\arraystretch}{1.25}
\begin{tabular}{@{}l@{\hspace{8mm}}ccccc@{}}
$f(x)$ & $\formule{x^2}$ & $\formule{x^3}$ & $\formule{\sqrt{x}}$ & $\formule{|x|}$ & $\formule{\frac1x}$ \\
$\dom f$ & $\R$ & $\R$ & $\left[0, +\infty\right[$ & $\R$ & $\R \setminus \{0\}$ \\
$\ber f$ & $\left[0, +\infty\right[$ & $\R$ & $\left[0, +\infty\right[$ & $\left[0, +\infty\right[$ & $\R \setminus \{0\}$ \\
symmetrie & even & oneven & geen & even & oneven
\end{tabular}
\end{theorie}

Je hoeft de tabel niet blind te leren: elke kolom volgt uit het voorschrift, met de begrippen uit hoofdstuk 1. Voor het bereik zeg je telkens welke waarden mogelijk zijn, en dat elke zulke waarde ook bereikt wordt.

\begin{itemize}
\item $x^2$: een kwadraat is nooit negatief, en elke $y \geqslant 0$ is het kwadraat van $\sqrt{y}$. Even, want $(-x)^2 = x^2$.
\item $x^3$: elke $y$ is de derdemacht van $\sqrt[3]{y}$, dus het bereik is $\R$. Oneven, want $(-x)^3 = -x^3$.
\item $\sqrt{x}$: bestaat alleen voor $x \geqslant 0$ en is nooit negatief; elke $y \geqslant 0$ is $\sqrt{y^2}$. Het domein is niet symmetrisch, dus geen symmetrie.
\item $|x|$: nooit negatief, en $|y| = y$ voor $y \geqslant 0$. Even, want $|-x| = |x|$.
\item $\frac1x$: bestaat niet in $0$, en is nooit $0$; elke $y \neq 0$ is het beeld van $\frac1y$. Oneven, want $\frac{1}{-x} = -\frac1x$.
\end{itemize}

\begin{oefeningen}
\oef{h2-basis}
\end{oefeningen}

%========================================================================================
\section{Verschuiven}

Een bakker zet elke ochtend zijn oven aan. Noem $f(t)$ de temperatuur van zijn oven om $t$ uur. Een collega doet precies hetzelfde, maar 2 uur vroeger. Haar oven heeft om $t$ uur de temperatuur die de oven van de bakker om $t + 2$ uur heeft: haar temperatuur is $f(t + 2)$. Is de oven van de bakker om 8 uur het warmst, dan is die van haar het warmst om 6 uur, want $f(6 + 2) = f(8)$. Er staat $+2$ in het voorschrift, maar het warmste moment komt \textit{vroeger}: de grafiek schuift naar links. Ook in de Verken was het zo: $g(x) = (x - 1)^2$ is de parabool $y = x^2$, één eenheid naar \textit{rechts} geschoven. Een verandering binnen de functie lijkt dus omgekeerd te werken. In deze sectie leggen we uit waarom. We bekijken eerst een verschuiving in de hoogte, dan een opzij, en ten slotte beide samen.

\subsection{Verticaal verschuiven}

Wat gebeurt er met een grafiek als je bij elke functiewaarde hetzelfde getal $d$ optelt? Dan komt elk punt van de grafiek $d$ hoger te liggen.

\begin{theorie}
\begin{stelling}{Verticale verschuiving}
De grafiek van $\formule{g(x) = f(x) + d}$ is die van $f$, verschoven over $d$ naar boven (naar onder als $d < 0$): \quad $\formule{(x, y) \mapsto (x, y + d)}$
\end{stelling}

\begin{alleenvoor}{3,4}
\begin{bewijs}
$g$ en $f$ hebben hetzelfde domein. Een punt $(x, y)$ ligt op de grafiek van $g$ als en slechts als
\bewijsregel{y = f(x) + d \iff y - d = f(x) ,}
dus als en slechts als $(x, y - d)$ op de grafiek van $f$ ligt. De punten van $g$ zijn dus precies de punten van $f$ met $d$ opgeteld bij de $y$-coördinaat.
\end{bewijs}
\end{alleenvoor}
\end{theorie}

Het domein verandert niet, het bereik schuift $d$ mee op.

\voorbeeld{}{Schets de grafiek van $g(x) = |x| - 2$, en geef het bereik en de nulwaarden.\label{vb:h2-vertverschuiven}}{%
\naast{%
\uitleg{Dit is $|x|$ met $d = -2$. Volg de knik $(0, 0)$.}
\begin{stappenvrij}{4.2cm}{4.6cm}
(0, 0) \mapsto (0, -2) & De knik zakt 2 eenheden. \\
\ber g = \left[-2, +\infty\right[ & Het bereik $\left[0, +\infty\right[$ zakt mee. \\
|x| - 2 = 0 \iff x = \pm 2 & De nulwaarden.
\end{stappenvrij}}{%
\begin{tikzpicture}[x=.45cm,y=.45cm,font=\scriptsize]
  \rooster{-3}{1}{3}{-2}{1}{2}
  \draw[grijs,dashed,thick] (-2.4,2.4) -- (0,0) -- (2.4,2.4);
  \draw[accent,thick] (-3,1) -- (0,-2) -- (3,1);
  \vol{(0,-2)}
\end{tikzpicture}}}

\begin{oefeningen}
\oef{h2-vert-kort}
\end{oefeningen}

\subsection{Horizontaal verschuiven}

Bij $g(x) = f(x - c)$ verandert iets aan de $x$. Wat $f$ doet in $0$, doet $g$ in $x = c$, want $g(c) = f(0)$. Alles gebeurt dus $c$ eenheden later.

\begin{theorie}
\begin{stelling}{Horizontale verschuiving}
De grafiek van $\formule{g(x) = f(x - c)}$ is die van $f$, verschoven over $c$ naar rechts (naar links als $c < 0$): \quad $\formule{(x, y) \mapsto (x + c, y)}$
\end{stelling}

\begin{alleenvoor}{3,4}
\begin{bewijs}
Een punt $(x, y)$ ligt op de grafiek van $g$ als en slechts als
\bewijsregel{y = f(x - c) ,}
dus als en slechts als $(x - c, y)$ op de grafiek van $f$ ligt. Noem $u = x - c$. Bij elk punt $(u, y)$ van $f$ hoort dus het punt $(u + c, y)$ van $g$, en omgekeerd.
\end{bewijs}
\end{alleenvoor}
\end{theorie}

Let op het teken: in $f(x - 3)$ schuift de grafiek naar \textit{rechts}, in $f(x + 3) = f\bigl(x - (-3)\bigr)$ naar links. Zo is de temperatuur van de oven van de collega $f(t + 2)$: $c = -2$, en de grafiek schuift 2 uur naar links. Nu verandert het domein mee, en blijft het bereik hetzelfde.

\voorbeeld{}{Schets de grafiek van $g(x) = \dfrac{1}{x - 1}$, en geef het domein.\label{vb:h2-horverschuiven}}{%
\naast{%
\uitleg{Dit is $\frac1x$ met $c = 1$. Volg twee punten, en de verticale rechte waar de functie niet bestaat.}
\begin{stappenvrij}{4.2cm}{4.6cm}
(1, 1) \mapsto (2, 1) & 1 naar rechts. \\
(-1, -1) \mapsto (0, -1) & Controle: $g(0) = -1$. \\
x = 0 \ \text{wordt}\ x = 1 & Daar bestaat $g$ niet. \\
\dom g = \R \setminus \{1\} & Het bereik blijft $\R \setminus \{0\}$.
\end{stappenvrij}}{%
\begin{tikzpicture}[x=.45cm,y=.45cm,font=\scriptsize]
  \rooster{-2}{1}{4}{-3}{1}{3}
  \draw[grijs,dashed] (1,-3) -- (1,3.3);
  \kromme{1.34}{4}{1/(\x-1)} \kromme{-2}{.66}{1/(\x-1)}
  \vol{(2,1)} \vol{(0,-1)}
\end{tikzpicture}}}

\begin{oefeningen}
\oef{h2-hor-kort}
\end{oefeningen}

\subsection{Beide samen}

Bij $g(x) = f(x - c) + d$ pas je beide stellingen na elkaar toe: eerst $c$ opzij, dan $d$ omhoog. Een punt $(x, y)$ gaat zo naar $(x + c, y + d)$. De volgorde maakt hier niet uit, want de ene verschuiving verandert alleen de $x$, de andere alleen de $y$.

\voorbeeld{}{Schets de grafiek van $g(x) = \sqrt{x + 3} - 1$, en geef het domein en het bereik.\label{vb:h2-verschuiven}}{%
\naast{%
\uitleg{Schrijf $x + 3$ als $x - (-3)$: dat is $c = -3$, en $d = -1$. Volg het beginpunt $(0, 0)$ van $\sqrt{x}$.}
\begin{stappenvrij}{4.3cm}{3.3cm}
c = -3,\ d = -1 & 3 naar links, 1 naar onder. \\
(0, 0) \mapsto (-3, -1) & Het beginpunt. \\
\dom g = \left[-3, +\infty\right[ & Ook 3 naar links. \\
\ber g = \left[-1, +\infty\right[ & Ook 1 naar onder.
\end{stappenvrij}}{%
\begin{tikzpicture}[x=.5cm,y=.5cm,font=\scriptsize]
  \rooster{-3}{1}{4}{-1}{1}{3}
  \draw[grijs,dashed,thick] plot[domain=0:4,samples=60,smooth] (\x,{sqrt(\x)});
  \kromme{-3}{4}{sqrt(\x+3)-1}
  \vol{(-3,-1)}
\end{tikzpicture}}}

\nujij{Waar ligt de top van de parabool $y = (x - 2)^2 + 1$?}{De top $(0, 0)$ van $y = x^2$ schuift 2 naar rechts en 1 naar boven: $(2, 1)$.}

\begin{oefeningen}
\oef{h2-verschuiven}

\needspace{8\baselineskip}
\oef{h2-verschuiven-grafiek}
\end{oefeningen}

%========================================================================================
\section{Uitrekken en spiegelen}

Een verschuiving telt iets op. Wat gebeurt er met de grafiek als je vermenigvuldigt, bijvoorbeeld elke functiewaarde met $2$? Vermenigvuldig je met een getal, dan wordt de grafiek uitgerekt of samengedrukt. Ook hier hangt het ervan af of je dat buiten $f$ doet of binnen.

\subsection{Verticaal uitrekken}

\begin{theorie}
\begin{stelling}{Verticaal uitrekken}
Voor $a > 0$ is de grafiek van $\formule{g(x) = a \cdot f(x)}$ die van $f$, verticaal uitgerekt met factor $a$: \quad $\formule{(x, y) \mapsto (x, ay)}$
\end{stelling}

\begin{alleenvoor}{3,4}
\begin{bewijs}
$g$ en $f$ hebben hetzelfde domein. Een punt $(x, y)$ ligt op de grafiek van $g$ als en slechts als
\bewijsregel{y = a f(x) \iff \frac{y}{a} = f(x) ,}
want $a \neq 0$. Dat is: als en slechts als $\bigl(x, \frac{y}{a}\bigr)$ op de grafiek van $f$ ligt. Bij elk punt $(x, v)$ van $f$ hoort dus het punt $(x, av)$ van $g$, en omgekeerd.
\end{bewijs}
\end{alleenvoor}
\end{theorie}

Is $a > 1$, dan wordt de grafiek steiler; is $0 < a < 1$, dan wordt ze platter. Punten op de $x$-as blijven liggen, want $a \cdot 0 = 0$: de nulwaarden veranderen niet.

\voorbeeld{}{Schets de grafiek van $g(x) = \frac12 x^2$.\label{vb:h2-vertrekken}}{%
\naast{%
\uitleg{Dit is $x^2$ met $a = \frac12$: elke hoogte halveert. Volg enkele punten van $x^2$.}
\begin{stappenvrij}{4.2cm}{4.4cm}
(0, 0) \mapsto (0, 0) & De top blijft liggen. \\
(\pm 2, 4) \mapsto (\pm 2, 2) & Half zo hoog. \\
(\pm 3, 9) \mapsto (\pm 3, \tfrac92) & Controle: $\frac12 \cdot 9 = \frac92$.
\end{stappenvrij}
De parabool is breder geworden, maar ze is niet opzij geschoven.}{%
\begin{tikzpicture}[x=.45cm,y=.4cm,font=\scriptsize]
  \rooster{-3}{1}{3}{0}{1}{5}
  \draw[grijs,dashed,thick] plot[domain=-2.2:2.2,samples=60,smooth] (\x,{(\x)^2});
  \kromme{-3.1}{3.1}{(\x)^2/2}
  \vol{(-2,2)} \vol{(2,2)}
\end{tikzpicture}}}

\begin{oefeningen}
\oef{h2-vertrekken-kort}
\end{oefeningen}

\subsection{Horizontaal uitrekken}

\begin{theorie}
\begin{stelling}{Horizontaal uitrekken}
Voor $b > 0$ is de grafiek van $\formule{h(x) = f(bx)}$ die van $f$, horizontaal uitgerekt met factor $\frac1b$: \quad $\formule{(x, y) \mapsto \left(\frac{x}{b}, y\right)}$
\end{stelling}

\begin{alleenvoor}{3,4}
\begin{bewijs}
Een punt $(x, y)$ ligt op de grafiek van $h$ als en slechts als
\bewijsregel{y = f(bx) ,}
dus als en slechts als $(bx, y)$ op de grafiek van $f$ ligt. Noem $u = bx$, dus $x = \frac{u}{b}$. Bij elk punt $(u, y)$ van $f$ hoort het punt $\bigl(\frac{u}{b}, y\bigr)$ van $h$, en omgekeerd.
\end{bewijs}
\end{alleenvoor}
\end{theorie}

Ook hier werkt een factor binnen de haakjes \textit{omgekeerd}: $f(2x)$ maakt de grafiek twee keer smaller, want $h(1) = f(2)$: wat $f$ in $2$ doet, doet $h$ al in $1$. Punten op de $y$-as blijven liggen, want $\frac0b = 0$.

\voorbeeld{}{De parabool $f(x) = x^2 - 4x$ heeft de nulwaarden $0$ en $4$ en de top $(2, -4)$. Bepaal de nulwaarden en de top van $h(x) = f(2x)$, en controleer met het voorschrift.\label{vb:h2-horrekken}}{%
\naast{%
\uitleg{Met $b = 2$ halveren alle $x$-coördinaten, de $y$-coördinaten blijven.}
\begin{stappenvrij}{4.2cm}{4.4cm}
0 \mapsto 0,\ 4 \mapsto 2 & Nulwaarden van $h$: $0$ en $2$. \\
(2, -4) \mapsto (1, -4) & De top van $h$. \\
h(x) = (2x)^2 - 4 \cdot 2x & Controle met het voorschrift: \\
\phantom{h(x)} = 4x(x - 2) & nulwaarden $0$ en $2$, \\
h(1) = 4 - 8 = -4 & en de top klopt.
\end{stappenvrij}}{%
\begin{tikzpicture}[x=.45cm,y=.4cm,font=\scriptsize]
  \rooster{-1}{1}{5}{-4}{1}{2}
  \draw[grijs,dashed,thick] plot[domain=-.45:4.45,samples=60,smooth] (\x,{(\x)^2-4*\x});
  \kromme{-.4}{2.4}{4*(\x)^2-8*\x}
  \vol{(1,-4)} \fill[grijs] (2,-4) circle (2.2pt);
\end{tikzpicture}}}

\begin{oefeningen}
\oef{h2-horrekken-kort}
\end{oefeningen}

\subsection{Verticaal of horizontaal?}

Het verschil tussen $a \cdot f(x)$ en $f(bx)$ zie je het best naast elkaar, met dezelfde grafiek en dezelfde factor.

\voorbeeld{}{De grafiek van $f$ is de gebroken lijn door $(-2, 0)$, $(0, 2)$, $(2, 0)$ en $(4, -2)$. Teken de grafieken van $g(x) = 2f(x)$ en $h(x) = f(2x)$.\label{vb:h2-rekken}}{%
\uitleg{Transformeer de vier hoekpunten en verbind ze. Bij $g$ verdubbelen de $y$-waarden, bij $h$ halveren de $x$-waarden.}
\begin{stappenbreed}
g\colon\ (-2, 0),\ (0, 4),\ (2, 0),\ (4, -4) & $(x, y) \mapsto (x, 2y)$: twee keer zo hoog. \\
h\colon\ (-1, 0),\ (0, 2),\ (1, 0),\ (2, -2) & $(x, y) \mapsto (\frac{x}{2}, y)$: twee keer zo smal.
\end{stappenbreed}
\figuren{%
\begin{tikzpicture}[x=.42cm,y=.42cm,font=\scriptsize]
  \rooster{-2}{1}{4}{-4}{1}{4}
  \draw[grijs,dashed,thick] (-2,0) -- (0,2) -- (2,0) -- (4,-2);
  \draw[accent,thick] (-2,0) -- (0,4) -- (2,0) -- (4,-4);
  \node[right,text=accent] at (.4,3.6) {$g$};
\end{tikzpicture}\hspace{1cm}
\begin{tikzpicture}[x=.42cm,y=.42cm,font=\scriptsize]
  \rooster{-2}{1}{4}{-4}{1}{4}
  \draw[grijs,dashed,thick] (-2,0) -- (0,2) -- (2,0) -- (4,-2);
  \draw[accent,thick] (-1,0) -- (0,2) -- (1,0) -- (2,-2);
  \node[right,text=accent] at (1.6,-2.3) {$h$};
\end{tikzpicture}}
Bij $g$ blijven de nulwaarden liggen, bij $h$ het snijpunt $(0, 2)$ met de $y$-as.}

\subsection{Spiegelen om de $x$-as}

Wat gebeurt er als de factor negatief is? Neem eerst $a = -1$: elke functiewaarde verandert van teken.

\begin{theorie}
\begin{stelling}{Spiegelen om de $x$-as}
De grafiek van $\formule{g(x) = -f(x)}$ is het spiegelbeeld van die van $f$ om de $x$-as: \quad $\formule{(x, y) \mapsto (x, -y)}$
\end{stelling}

\begin{alleenvoor}{3,4}
\begin{bewijs}
Het bewijs van verticaal uitrekken gebruikt alleen dat $a \neq 0$. Met $a = -1$ geeft het: bij elk punt $(x, v)$ van $f$ hoort het punt $(x, -v)$ van $g$. Dat is het spiegelbeeld om de $x$-as.
\end{bewijs}
\end{alleenvoor}
\end{theorie}

Een negatieve factor combineert dus twee transformaties: $-3f(x) = -\bigl(3f(x)\bigr)$ is verticaal uitrekken met factor $3$ en daarna spiegelen om de $x$-as.

\voorbeeld{}{Geef het domein en het bereik van $k(x) = -\sqrt{x}$.\label{vb:h2-spiegelx}}{%
\naast{%
\uitleg{Spiegel $\sqrt{x}$ om de $x$-as: de $x$-coördinaten blijven, de $y$-coördinaten veranderen van teken.}
\begin{stappenvrij}{4.2cm}{4.2cm}
\dom k = \left[0, +\infty\right[ & Zoals bij $\sqrt{x}$. \\
\ber k = \left]-\infty, 0\right] & Het bereik klapt om.
\end{stappenvrij}}{%
\begin{tikzpicture}[x=.45cm,y=.45cm,font=\scriptsize]
  \rooster{-1}{1}{4}{-2}{1}{2}
  \draw[grijs,dashed,thick] plot[domain=0:4,samples=60,smooth] (\x,{sqrt(\x)});
  \kromme{0}{4}{-sqrt(\x)}
  \vol{(0,0)}
\end{tikzpicture}}}

\subsection{Spiegelen om de $y$-as}

Met $b = -1$ verandert binnen de haakjes $x$ in $-x$: wat $f$ in $2$ doet, doet de nieuwe functie in $-2$.

\begin{theorie}
\begin{stelling}{Spiegelen om de $y$-as}
De grafiek van $\formule{h(x) = f(-x)}$ is het spiegelbeeld van die van $f$ om de $y$-as: \quad $\formule{(x, y) \mapsto (-x, y)}$
\end{stelling}

\begin{alleenvoor}{3,4}
\begin{bewijs}
Het bewijs van horizontaal uitrekken gebruikt alleen dat $b \neq 0$. Met $b = -1$ geeft het: bij elk punt $(u, y)$ van $f$ hoort het punt $(-u, y)$ van $h$. Dat is het spiegelbeeld om de $y$-as.
\end{bewijs}
\end{alleenvoor}
\end{theorie}

Nu klapt het domein om, en blijft het bereik. Dit sluit aan bij hoofdstuk 1: voor een even functie is $f(-x) = f(x)$, en spiegelen om de $y$-as verandert de grafiek niet. Voor een oneven functie is $f(-x) = -f(x)$: beide spiegelingen geven dan hetzelfde.

\voorbeeld{}{Geef het domein en het bereik van $h(x) = \sqrt{-x}$.\label{vb:h2-spiegely}}{%
\naast{%
\uitleg{Spiegel $\sqrt{x}$ om de $y$-as: de $y$-coördinaten blijven, de $x$-coördinaten veranderen van teken.}
\begin{stappenvrij}{4.2cm}{4.2cm}
\dom h = \left]-\infty, 0\right] & Het domein klapt om. \\
\ber h = \left[0, +\infty\right[ & Zoals bij $\sqrt{x}$.
\end{stappenvrij}
Controle: $h(-4) = \sqrt{4} = 2$, en $(-4, 2)$ is het spiegelbeeld van $(4, 2)$.}{%
\begin{tikzpicture}[x=.45cm,y=.45cm,font=\scriptsize]
  \rooster{-4}{1}{4}{-1}{1}{2}
  \draw[grijs,dashed,thick] plot[domain=0:4,samples=60,smooth] (\x,{sqrt(\x)});
  \kromme{-4}{0}{sqrt(-\x)}
  \vol{(0,0)}
\end{tikzpicture}}}

\nujij{Welke functie krijg je als je de grafiek van $x^3$ om de $y$-as spiegelt? En om de $x$-as?}{Om de $y$-as: $(-x)^3 = -x^3$. Om de $x$-as: $-x^3$. Hetzelfde, want $x^3$ is oneven.}

\begin{oefeningen}
\oef{h2-rekken}

\oef{h2-spiegelen}
\end{oefeningen}

%========================================================================================
\section{Transformaties combineren}

Bij een voorschrift als $g(x) = a \cdot f\bigl(b(x - c)\bigr) + d$ voer je de transformaties na elkaar uit. Welke volgorde is de goede? Volg een punt van $f$ en kijk waar het terechtkomt.

\begin{theorie}
\begin{stelling}{Een punt volgen}
Is $\formule{g(x) = a \cdot f\bigl(b(x - c)\bigr) + d}$ met $a, b \neq 0$, dan hoort bij elk punt $(u, v)$ van de grafiek van $f$ het punt \ $\formule{\bigl(\tfrac{u}{b} + c,\ av + d\bigr)}$ \ van de grafiek van $g$, en omgekeerd.
\end{stelling}

\begin{alleenvoor}{3,4}
\begin{bewijs}
Neem een punt $(x, y)$ en noem $u = b(x - c)$ en $v = \frac{y - d}{a}$, zodat $x = \frac{u}{b} + c$ en $y = av + d$. Dan is
\bewijsregel{y = a f\bigl(b(x - c)\bigr) + d \iff av + d = a f(u) + d \iff v = f(u) .}
Het punt $(x, y)$ ligt dus op de grafiek van $g$ als en slechts als $(u, v)$ op die van $f$ ligt.
\end{bewijs}
\end{alleenvoor}
\end{theorie}

Lees de formule als een recept. Voor de $x$-coördinaat deel je eerst door $b$ en tel je daarna $c$ op: eerst horizontaal uitrekken, dan verschuiven. Voor de $y$-coördinaat vermenigvuldig je eerst met $a$ en tel je daarna $d$ op: eerst verticaal uitrekken, dan verschuiven. Werk dus van binnen naar buiten.

\begin{alleenvoor}{1,2,3}
Met ICT zie je wat elke parameter doet. Maak in GeoGebra vier schuifbalken $a$, $b$, $c$ en $d$, en voer in: \texttt{g(x) = a*f(b*(x-c))+d}, met bijvoorbeeld \texttt{f(x) = sqrt(x)}. Verschuif telkens één schuifbalk, en volg wat er met het beginpunt gebeurt.

\end{alleenvoor}
\voorbeeld{}{Schets $g(x) = -2\sqrt{x + 1} + 3$ vanuit $\sqrt{x}$, en geef het domein en het bereik.\label{vb:h2-combineren}}{%
\naast{%
\uitleg{Volg het beginpunt $(0, 0)$ door alle stappen.}
\begin{stappenvrij}{3.9cm}{3.7cm}
\sqrt{x + 1} & 1 naar links: $(-1, 0)$. \\
2\sqrt{x + 1} & Verticaal $\times 2$: $(-1, 0)$. \\
-2\sqrt{x + 1} & Spiegelen om de $x$-as. \\
-2\sqrt{x + 1} + 3 & 3 naar boven: $(-1, 3)$. \\
\dom g = \left[-1, +\infty\right[ & \\
\ber g = \left]-\infty, 3\right] & De grafiek daalt vanaf $3$.
\end{stappenvrij}}{%
\begin{tikzpicture}[x=.5cm,y=.5cm,font=\scriptsize]
  \rooster{-1}{1}{5}{-2}{1}{3}
  \draw[grijs,dashed,thick] plot[domain=0:5,samples=60,smooth] (\x,{sqrt(\x)});
  \kromme{-1}{5}{-2*sqrt(\x+1)+3}
  \vol{(-1,3)}
\end{tikzpicture}}}

\begin{werkwijze}
Schrijf het voorschrift als $a \cdot f\bigl(b(x - c)\bigr) + d$. Pas toe in deze volgorde: horizontaal uitrekken met $\frac1b$ (en spiegelen als $b < 0$), verschuiven over $c$ naar rechts, verticaal uitrekken met $a$ (en spiegelen als $a < 0$), verschuiven over $d$ naar boven. Volg een bijzonder punt, en controleer met één functiewaarde.
\end{werkwijze}

\begin{oefeningen}
\oef{h2-combineren-schets}
\end{oefeningen}

Staat er een factor bij $x$, zet die dan eerst buiten haakjes. Anders lees je de verschuiving verkeerd af.

\voorbeeld{}{Schets $g(x) = \sqrt{2x - 6}$ vanuit $\sqrt{x}$.\label{vb:h2-buitenhaakjes}}{%
\naast{%
\uitleg{Een vlugge lezing zegt: 6 naar rechts. Maar dan zou het beginpunt $(6, 0)$ zijn, terwijl $g(6) = \sqrt6 \neq 0$. Zet eerst $2$ buiten haakjes.}
\begin{stappenvrij}{4.2cm}{4.4cm}
g(x) = \sqrt{2(x - 3)} & $b = 2$ en $c = 3$. \\
(0, 0) \mapsto (3, 0) & $\frac02 + 3 = 3$. \\
(4, 2) \mapsto (5, 2) & $\frac42 + 3 = 5$. \\
g(5) = \sqrt{4} = 2 & Controle.
\end{stappenvrij}
Eerst horizontaal $\times \frac12$, dan $3$ (niet $6$) naar rechts.}{%
\begin{tikzpicture}[x=.45cm,y=.45cm,font=\scriptsize]
  \rooster{-1}{1}{6}{-1}{1}{3}
  \draw[grijs,dashed,thick] plot[domain=0:6,samples=60,smooth] (\x,{sqrt(\x)});
  \kromme{3}{6}{sqrt(2*\x-6)}
  \vol{(3,0)} \vol{(5,2)}
\end{tikzpicture}}}

\begin{oefeningen}
\oef{h2-combineren}
\end{oefeningen}

Omgekeerd kun je uit een grafiek het voorschrift halen, als je weet van welke basisfunctie ze komt. Een bijzonder punt geeft $c$ en $d$, een tweede punt de factor $a$.

\voorbeeld{}{De parabool hiernaast heeft top $(-1, 2)$ en gaat door $(1, 0)$. Geef het voorschrift.\label{vb:h2-voorschrift}}{%
\naast{%
\uitleg{Een parabool is $x^2$, verschoven naar de top en verticaal uitgerekt: $y = a(x - c)^2 + d$. Bepaal $a$ met het tweede punt.}
\begin{stappenvrij}{4.6cm}{3.6cm}
y = a(x + 1)^2 + 2 & Top $(-1, 2)$. \\
0 = a \cdot 2^2 + 2 \Rightarrow a = -\tfrac12 & Punt $(1, 0)$ invullen. \\
y = -\tfrac12(x + 1)^2 + 2 &
\end{stappenvrij}}{%
\begin{tikzpicture}[x=.5cm,y=.5cm,font=\scriptsize]
  \rooster{-4}{1}{2}{-1}{1}{3}
  \kromme{-3.5}{1.5}{-(\x+1)^2/2+2}
  \vol{(-1,2)} \vol{(1,0)}
\end{tikzpicture}}}

\begin{oefeningen}
\oef{h2-voorschrift-kort}
\end{oefeningen}

\voorbeeld{}{Terug naar het probleem. $f(t)$ is de temperatuur in Brussel op een winterdag, in $^\circ$C, $t$ uur na middernacht (Brusselse tijd), met domein $[0, 24]$. De top van de grafiek is $(15, 6)$, het dal $(3, -2)$, en de nulwaarden zijn $7$ en $23$. Je vriendin in New York, waar het 6 uur vroeger is dan in Brussel, wil dezelfde metingen zien zoals ze dat gewend is: in graden Fahrenheit en volgens de klok van New York. Een temperatuur van $T$\,$^\circ$C is $1{,}8\,T + 32$\,$^\circ$F. Teken haar grafiek. Om hoe laat was het volgens haar klok het warmst, en hoeveel graden Fahrenheit was het toen? Waar zie je op haar grafiek dat het vroor?\label{vb:h2-terug}}{%
\uitleg{Om $t$ uur in New York is het $t + 6$ uur in Brussel, en de temperatuur $f(t + 6)$ in $^\circ$C. Reken die om naar $^\circ$F. Schrijf het voorschrift in de vorm van de stelling \textit{Een punt volgen}, en volg de bijzondere punten. Zo hoef je geen enkele temperatuur opnieuw uit te rekenen.}
\begin{stappenbreed}
g(t) = 1{,}8 \cdot f\bigl(t - (-6)\bigr) + 32 & $a = 1{,}8$, $b = 1$, $c = -6$, $d = 32$. \\
(u, v) \mapsto (u - 6;\ 1{,}8v + 32) & Elk punt: 6 uur naar links, de hoogte maal $1{,}8$ en dan $32$ hoger. \\
(15, 6) \mapsto (9;\ 42{,}8) & Het warmst: om 9 uur, $42{,}8$\,$^\circ$F. \\
(3, -2) \mapsto (-3;\ 28{,}4) & Het koudst: om $-3$ uur, dat is 21 uur de avond ervoor. \\
(7, 0) \mapsto (1, 32),\ \ (23, 0) \mapsto (17, 32) & De nulwaarden komen op de rechte $T = 32$. \\
\dom g = [-6, 18] & Het domein schuift mee: van 18 uur de dag ervoor tot 18 uur.
\end{stappenbreed}
\naast{%
Het vriespunt $0$\,$^\circ$C is $32$\,$^\circ$F: de $t$-as van Brussel wordt de rechte $T = 32$ (gestippeld). Het vroor dus waar de grafiek onder die rechte ligt: volgens de klok van New York tot 1 uur 's nachts en vanaf 17 uur. Controle: om 9 uur in New York is het 15 uur in Brussel, en $1{,}8 \cdot 6 + 32 = 42{,}8$. De grafiek heeft dezelfde vorm als die van Brussel: alleen de assen kregen andere getallen.}{%
\begin{tikzpicture}[x=.22cm,y=.13cm,font=\scriptsize]
  \draw[black!14,very thin,xstep=.66cm,ystep=.52cm] (-6,24) grid (18,44);
  \draw[->,grijs] (-6,24) -- (19.5,24) node[right,text=black] {$t$};
  \draw[->,grijs] (-6,24) -- (-6,46.5) node[above,text=black] {$^\circ$F};
  \foreach \t in {-6,0,6,12,18}{\node[below,inner sep=1.5pt] at (\t,24) {$\t$};}
  \foreach \t in {28,32,36,40,44}{\node[left,inner sep=1.5pt] at (-6,\t) {$\t$};}
  \draw[grijs,dashed] (-6,32) -- (18,32);
  \kromme{-6}{18}{-7.2*cos(15*(\x+3))+35.6}
  \vol{(9,42.8)} \vol{(-3,28.4)} \vol{(1,32)} \vol{(17,32)}
\end{tikzpicture}}}

%========================================================================================
\opdrachtenkop

\begin{oefeningen}
\needspace{10\baselineskip}
\oef{h2-o1}

\needspace{9\baselineskip}
\oef{h2-o10}

\oef{h2-o2}

\needspace{9\baselineskip}
\oef{h2-o11}

\needspace{7\baselineskip}
\oef{h2-o14}

\needspace{6\baselineskip}
\oef{h2-o12}

\begin{alleenvoor}{1,2,3}
\needspace{12\baselineskip}
\oef{h2-schuifbalken}

\end{alleenvoor}
\needspace{8\baselineskip}
\oef{h2-o13}

\needspace{10\baselineskip}
\oef{h2-o16}

\needspace{9\baselineskip}
\oef{h2-o3}

\needspace{7\baselineskip}
\oef{h2-o20}

\needspace{8\baselineskip}
\oef{h2-o4}

\needspace{10\baselineskip}
\oef{h2-o6}

\needspace{9\baselineskip}
\oef{h2-o17}

\needspace{11\baselineskip}
\oef{h2-o5}

\needspace{9\baselineskip}
\oef{h2-o21}

\needspace{8\baselineskip}
\oef{h2-o24}

\needspace{9\baselineskip}
\oef{h2-o18}

\needspace{10\baselineskip}
\oef{h2-o25}

\needspace{10\baselineskip}
\oef{h2-o7}

\needspace{10\baselineskip}
\oef{h2-o27}

\needspace{8\baselineskip}
\oef{h2-o26}

\oef{h2-o8}

\needspace{8\baselineskip}
\oef{h2-o28}

\oef{h2-o9}
\end{oefeningen}

%----------------------------------------------------------------------------------------
% KOPPEL: verbanden met andere leerstof
%----------------------------------------------------------------------------------------
\needspace{\alleenvoor{2,3,4}{20}\alleenvoor{1}{8}\baselineskip}
\opdrachtsectie{Koppel}

\lkr{Goniometrische functies\alleenvoor{2,3,4}{, Exponentiële en logaritmische functies en Statistiek}\alleenvoor{1}{ en Exponentiële en logaritmische functies}: kijk na wat al gezien is. Het pakket legt de volgorde niet vast.}

\begin{oefeningen}
\oef{h2-koppel-1}

\oef{h2-koppel-2}

\oef{h2-koppel-3}
\end{oefeningen}

%----------------------------------------------------------------------------------------
\opdrachtsectie{Test jezelf}

\begin{oefeningen}
\oef{h2-t5}

\oef{h2-t1}

\oef{h2-t2}

\oef{h2-t6}

\oef{h2-t3}

\oef{h2-t7}

\oef{h2-t4}
\end{oefeningen}

\herhaalschema{%
  \alleenvoor{2,3,4}{\fzv{oef:h2-t5}{oef:h2-basis}}%
  \fv{oef:h2-t1}{vb:h2-verschuiven}{oef:h2-verschuiven}%
  \fv{oef:h2-t2}{vb:h2-rekken}{oef:h2-rekken}%
  \fv{oef:h2-t6}{vb:h2-rekken}{oef:h2-o10}%
  \fv{oef:h2-t3}{vb:h2-combineren}{oef:h2-combineren}%
  \fv{oef:h2-t7}{vb:h2-combineren}{oef:h2-combineren-schets}%
  \fv{oef:h2-t4}{vb:h2-voorschrift}{oef:h2-voorschrift-kort}}

\ifzelfstudie
\section*{Maak je eigen samenvatting}
Sluit dit boek en schrijf in je schrift: de grafieken van de vijf basisfuncties met domein en bereik, wat $f(x - c) + d$, $a f(x)$, $f(bx)$, $-f(x)$ en $f(-x)$ met een grafiek doen, en in welke volgorde je combineert. Vergelijk daarna met de kaders.
\fi
